\section*{Задания}
\addcontentsline{toc}{section}{Задания}

\subsection*{1. Представление чисел с фиксированной и плавающей запятой в различных форматах}
\addcontentsline{toc}{subsection}{Представление чисел с фиксированной и плавающей запятой в различных форматах}

\begin{enumerate}
    \item Заданное число $A$ представить в виде двоично-кодированного
    десятичного числа:
    \begin{enumerate}
        \item в упакованном формате (BCD);
        \item в неупакованном формате (ASCII).
    \end{enumerate}

    \item Заданные числа $A$ и $-A$ представить в форме с фиксированной
    запятой.

    \item Заданные числа $A$ и $B$ представить в форме с плавающей
    запятой в формате Ф1.

    \item Заданные числа $A$ и $B$ представить в форме с плавающей
    запятой в формате Ф2.

    \item Заданные числа $A$ и $B$ представить в форме с плавающей
    запятой в формате Ф3.

    \item Найти значения чисел $Y$ и $Z$ по их заданным
    шестнадцатеричным представлениям $R$ и $S$ в форме с плавающей
    запятой в формате Ф1.

    \item Найти значения чисел $V$ и $W$ по их заданным
    шестнадцатеричным представлениям $R$ и $S$ в форме с плавающей
    запятой в формате Ф2.

    \item Найти значения чисел $T$ и $Q$ по их заданным
    шестнадцатеричным представлениям $R$ и $S$ в форме с плавающей
    запятой в формате Ф3.
\end{enumerate}

\noindent\textbf{Замечание.}
При выполнении пп.~3--5 задания для дробного числа $B$ в целях
увеличения точности его представления произвести симметричное
округление мантиссы.

\subsection*{Заданные числа}
\addcontentsline{toc}{subsection}{Заданные числа}

\begin{table}[htbp]
    \centering
    \caption{Заданные значения из варианта 14}
    \begin{tabular}{|c|c|c|c|}
		$A$ & $B$ & $R$ & $S$ \\
		\hline
		$1400$ & $0.05$ & $C1E08000$ & $3EA00000$ \\
		
    \end{tabular}
\end{table}


\subsection*{2. Сложение целых чисел}
\addcontentsline{toc}{subsection}{Сложение целых чисел}

\begin{enumerate}
    \item Для заданных чисел $A$ и $B$ выполнить операцию знакового
    сложения со всеми комбинациями знаков операндов. Для каждого
    примера:
    \begin{enumerate}
        \item проставить межразрядные переносы, возникающие при сложении;
        \item дать знаковую интерпретацию (Зн) операндов и результатов.
        При получении отрицательного результата предварительно
        преобразовать его из дополнительного кода в прямой;
        \item дать беззнаковую интерпретацию (БзИ) операндов и
        результатов. При получении неверного результата пояснить
        причину его возникновения;
        \item показать значения арифметических флагов.
    \end{enumerate}

    \item Сохранив значение первого операнда $A$, выбрать такое
    значение $B$, чтобы в операции сложения с одинаковыми знаками
    имел место особый случай переполнения формата. Выполнить два
    примера, иллюстрирующие эти случаи. Для каждого из них выполнить
    пункты а--d.

    \item Сохранив операнд $B$, подобрать такое значение операнда $A$,
    чтобы при сложении положительных операндов имело место
    переполнение формата, а при сложении отрицательных операндов
    результат был бы корректен. Выполнить два примера,
    иллюстрирующие этот случай. Для каждого из них выполнить
    пункты а--d.
\end{enumerate}

\subsection*{Заданные числа}
\addcontentsline{toc}{subsection}{Заданные числа}

\begin{table}[H]
    \centering
    \caption{Заданные значения из варианта 14}
    \begin{tabular}{|c|c|}
		$A$ & $B$ \\
		\hline
		$109$ & $17$  \\
		
    \end{tabular}
\end{table}



\subsection*{3. Вычитание целых чисел}
\addcontentsline{toc}{subsection}{Вычитание целых чисел}

\begin{enumerate}
    \item Для заданных чисел $A$ и $B$ выполнить операцию знакового
    вычитания со всеми комбинациями знаков операндов. Для каждого
    примера:
    \begin{enumerate}
        \item проставить межразрядные заёмы, возникающие при вычитании;
        \item дать знаковую интерпретацию (Зн) операндов и результатов.
        При получении отрицательного результата предварительно
        преобразовать его из дополнительного кода в прямой;
        \item дать беззнаковую интерпретацию (БзИ) операндов и
        результатов. При получении неверного результата пояснить
        причину его возникновения;
        \item показать значения арифметических флагов.
    \end{enumerate}

    \item Сохранив значение первого операнда $A$, выбрать такое
    значение $B$, чтобы в операции вычитания с разными знаками имел
    место особый случай переполнения формата. Выполнить два примера,
    иллюстрирующих эти случаи. Для каждого из них выполнить пункты а--d.

    \item Сохранив операнд $B$, подобрать такое значение операнда $A$,
    чтобы при вычитании отрицательного $B$ из положительного $A$ имело
    место переполнение формата, а при вычитании положительного $B$ из
    отрицательного $A$ результат был бы корректен. Выполнить два
    примера, иллюстрирующих этот случай. Для каждого из них выполнить
    пункты а--d.
\end{enumerate}

\subsection*{Заданные числа}
\addcontentsline{toc}{subsection}{Заданные числа}

\begin{table}[H]
    \centering
    \caption{Заданные значения из варианта 14}
    \begin{tabular}{|c|c|}
		$A$ & $B$ \\
		\hline
		$109$ & $17$  \\
		
    \end{tabular}
\end{table}

